%% AIAgNav --- ICRA 2027 submission
%% Build:  make          (pdf)
%%         make pages    (page count -- the only number that matters)
%%         make words    (per-section word count vs. budget)
%%
%% HARD LIMIT: 8 pages including references.
%% Budget measured from P-AgNav (RA-L 2025), which is exactly 8 pages:
%%
%%   Abstract + I. Intro .... 1344 w
%%   II. System Overview .....  989 w
%%   III. System Design ..... 2831 w   <-- the technical core
%%   IV. Experiments ........ 1003 w
%%   V. Conclusion ..........  247 w
%%   References .............  958 w
%%   TOTAL .................. 7372 w
%%
%% Figures are not overhead -- they are how P-AgNav stays short. Their Fig. 4
%% (pipeline diagram) replaces roughly a page of prose. Budget ~2 pages of
%% figures and write less text, not more.

\documentclass[conference]{IEEEtran}
\IEEEoverridecommandlockouts

\usepackage{graphicx}
\usepackage{amsmath,amssymb}
\usepackage{booktabs}
\usepackage{url}
\usepackage{cite}
\usepackage[caption=false,font=footnotesize]{subfig}

\graphicspath{{fig/}}

%% Unresolved facts render visibly in the PDF so they cannot slip through.
\newcommand{\TODO}[1]{\textbf{[TODO: #1]}}

%% --- notation shorthands: define once, use everywhere -------------------
\newcommand{\offset}{e_t}                 % normalized lateral offset
\newcommand{\slope}{s_t}                  % heading / slope term
\newcommand{\angvel}{\omega_t}            % angular velocity
\newcommand{\linvel}{v_t}                 % linear velocity
\newcommand{\horizon}{N}                  % MPC horizon

\begin{document}

\title{AIAgNav: A Semantic Segmentation-Based Autonomous\\
       Navigation System for Cornfields}

\author{\IEEEauthorblockN{Author One, Author Two, and David J. Cappelleri}
\IEEEauthorblockA{School of Mechanical Engineering, Purdue University,
West Lafayette, IN 47907 USA}%
\thanks{This work was supported by ...}}

\maketitle

%% ===================================================================
%% ABSTRACT  --- budget: ~150 w
%% ===================================================================
\begin{abstract}
TODO. One sentence of problem, one of approach, one of what is new,
one of validation. See P-AgNav's abstract for the exact shape.
\end{abstract}

%% ===================================================================
%% I. INTRODUCTION  --- budget: ~1200 w (with abstract: 1344 w)
%% Deferred per author's instruction; write last, once III is settled.
%% Ends with an explicit 3-item contribution list (P-AgNav does this).
%% ===================================================================
\section{Introduction}
TODO --- deferred.

%% ===================================================================
%% II. SYSTEM OVERVIEW  --- budget: ~990 w  (draft currently 2274 w)
%%   A. System Hardware
%%   B. Sensing modality: why a monocular camera + learned segmentation
%%   C. Framework Overview  (+ pipeline figure, + field assumptions)
%% This section motivates the representation. It is the AIAgNav analogue
%% of P-AgNav Sec. II-B ("LiDAR Representation: Range View") and should
%% carry the "why DINOv3 / why three classes" argument at the level of
%% a design justification, NOT a tutorial.
%% ===================================================================
\section{System Overview}
\subsection{System Hardware}
TODO
\subsection{Visual Representation}
TODO
\subsection{Framework Overview}
TODO

%% ===================================================================
%% III. SYSTEM DESIGN --- budget: 2831 w TOTAL (draft currently 11812 w)
%% Suggested split, mirroring P-AgNav's proportions:
%%   A. Semantic Segmentation .......... ~700 w
%%   B. Centerline / Row-Width Estimation ~600 w
%%   C. MPC Row Following .............. ~800 w   (equation-dense)
%%   D. Row Exit Detection ............. ~400 w
%%   E. Mission FSM / Row Switching .... ~330 w
%% Equations carry the content; prose is connective tissue only.
%% ===================================================================
\section{System Design}

\begin{table}[t]
\caption{Summary of Notations}
\label{tab:notation}
\centering
\begin{tabular}{ll}
\toprule
Symbol & Description \\
\midrule
$\offset$   & Normalized lateral offset of the row centerline \\
$\slope$    & Heading-free lateral readout (slope term) \\
$\angvel$   & Commanded angular velocity \\
$\linvel$   & Commanded linear velocity \\
$\horizon$  & MPC prediction horizon \\
\bottomrule
\end{tabular}
\end{table}

\subsection{Semantic Segmentation}
The navigation state is derived from a per-pixel semantic segmentation of a
single forward-facing RGB frame. Under a corn canopy the appearance cues that
separate drivable ground from crop are unstable: illumination changes by more
than an order of magnitude between direct sun and deep shade within one row,
hanging leaves throw high-contrast shadows across the corridor floor, and soil
colour shifts with moisture and tillage. Hand-tuned colour or texture
thresholds are adequate over a short segment but do not survive these
transitions, and the row-following controller carries no independent state
estimate to fall back on when they fail. The system therefore learns the
ground/crop distinction rather than specifying it.

The network pairs a DINOv3 ViT-S/16 backbone with an encoder-only mask
transformer head. The choice is driven by label efficiency rather than by
accuracy on a public benchmark. Annotated data for this task cannot be
inherited from an existing dataset: it must be produced by hand from footage
recorded in the target field, at the growth stage and camera mounting of the
deployed robot, and the achievable quantity is bounded by how much one operator
can label in a season. The training set here is \TODO{confirm image count}
annotated frames. Self-supervised DINOv3 features transfer to dense prediction
under very little supervision, which is what makes a set of that size viable;
a backbone trained from scratch, or one supervised on ImageNet classification,
would demand substantially more annotation before producing masks clean enough
to measure corridor boundaries on. ViT-S/16 is the smallest variant in the
family, which matters because one set of weights must serve both robots in the
fleet, only one of which carries a GPU.

The model predicts three classes: sky, traversable ground, and obstacle. The
controller reads only the traversable class, so the presence of a sky class
deserves an explanation. At a row end, and throughout a headland, the upper
part of the frame opens onto bright sky. A two-class model must assign those
pixels to either ground or crop, and their brightness and near-absence of
texture make them resemble open ground far more closely than they resemble a
corn plant. Because the exit detector of Section~\ref{sec:exit} fires when the
traversable corridor widens toward the full image width, folding sky into the
traversable class reproduces the exact signature of a row end at arbitrary
points inside a row. A dedicated sky label removes that failure mode for the
cost of one output channel. The obstacle class, conversely, is deliberately
coarse: it covers the corn plants bounding the corridor, fallen stalks, and
any foreign object, because the control decision they induce is identical --
do not drive there. Separating them would add annotation burden without
changing a single command.

Training runs for 2\,500 steps at batch size 2 and $224 \times 224$ input
resolution in mixed precision. The schedule is short by the standards of
segmentation training, and deliberately so: with a training set this small, a
longer schedule fits the particular plants that happened to be annotated rather
than the ground/crop boundary the controller needs.

At inference the model is given the full-resolution camera frame and returns a
class-index mask, which is resized to the frame resolution when the two differ.
That resize must be nearest-neighbour. Mask values are class indices, not
intensities, so any interpolating kernel synthesises intermediate values that
correspond to no class at all, and it does so precisely at region boundaries --
which is exactly where Section~\ref{sec:centerline} measures the corridor
edges. A single misplaced boundary column biases the lateral error for that
frame.

\subsection{Corridor and Centerline Estimation}
\label{sec:centerline}
The mask is reduced to a two-dimensional navigation state by pure geometry.
No regression head is trained to predict heading or lateral distance, and no
camera intrinsics enter the computation: the state is read directly off the
class mask, so it is defined in normalized image coordinates and is
independent of camera resolution.

The system measures the corridor on a small set of horizontal scan rows placed
at fractions $f_i$ of the image height, ordered from farthest to nearest. On
each scan row it starts at the image centre column $c_x = W/2$ and expands
outward while the mask remains traversable, giving the bounds $x_{L,i}$ and
$x_{R,i}$ of the contiguous traversable run that straddles the centre column,
\begin{equation}
\label{eq:midpoint}
x_{\mathrm{mid},i} = \tfrac{1}{2}\left(x_{L,i} + x_{R,i}\right),
\qquad
w_i = x_{R,i} - x_{L,i} .
\end{equation}
Anchoring the search at $c_x$ rather than taking the widest traversable run in
the row is what makes the measurement unambiguous when the mask contains more
than one open region, as it does whenever a gap in the crop wall exposes the
neighbouring row: the corridor the robot is actually in is the one aligned
with its own heading. If the centre column is not traversable, that scan row
yields no measurement and is skipped. The width $w_i$ is retained and consumed
by the exit detector of Section~\ref{sec:exit}.

Each valid scan row gives a lateral error, normalized by the half-width so
that it is dimensionless and bounded, and these are pooled into a single
lateral state $e$ by a weighted mean over the valid rows,
\begin{equation}
\label{eq:offset}
e_i = \frac{x_{\mathrm{mid},i} - c_x}{W/2},
\qquad
e = \mathrm{clip}\!\left(\frac{\sum_i \lambda_i e_i}{\sum_i \lambda_i},\, -1,\, 1\right).
\end{equation}
A positive $e$ places the corridor midpoint right of the image centre, meaning
the robot sits left of the corridor centre. The weights $\lambda_i$ favour the
near rows. Near rows and far rows fail in opposite ways, and the weighting
reflects which failure is more costly: the near rows report where the robot is
now and stay reliable because the ground closest to the camera segments
cleanly, whereas the far rows carry the earliest warning of a curve or a row
end but sit where perspective compresses the corridor to a few pixels and the
mask degrades. Weighting toward the near field keeps steering stable while
leaving the far rows enough influence to anticipate.

Because the same scan already yields an offset at several distances, the
robot's heading relative to the row follows from their difference at no
additional cost,
\begin{equation}
\label{eq:slope}
s = e_{\mathrm{far}} - e_{\mathrm{near}},
\end{equation}
evaluated between the farthest and nearest valid rows. A corridor that tilts
rightward with distance indicates the robot is heading left of the row
direction. This is the term that separates the two ways a robot can be wrong
inside a row -- displaced from the centreline, and pointed across it -- which
a single lateral reading cannot distinguish. It is computed across rows within
one frame rather than across time, so it needs no history, no filter, and no
inertial measurement.

A frame is accepted only if at least one scan row produced a measurement and
the traversable fraction $\tau$ of the lower half of the mask reaches a
threshold $\tau_{\min}$. The second test rejects the case where the mask has
collapsed -- an occluded lens, a lost row -- while a narrow spurious run at the
centre column still yields a confident-looking midpoint. On rejection the
controller holds its previous command rather than steering on the frame, and
stops if rejections persist.

\subsection{Model Predictive Row Following}
The MPC minimizes the cost $J$ over the horizon $\horizon$:
\begin{equation}
\label{eq:mpc_cost}
J = \sum_{k=1}^{\horizon} \big( \text{TODO} \big)
\end{equation}
subject to the state update and constraints
\begin{equation}
\label{eq:mpc_dynamics}
\text{TODO}
\end{equation}

\subsection{Row Exit Detection}
\label{sec:exit}
TODO

\subsection{Row Switching}
TODO

%% ===================================================================
%% IV. EXPERIMENTAL RESULTS --- budget: ~1000 w
%% Two tables (SIM, field) + the failure-mode figure. P-AgNav reports
%% collisions and human interventions; we report MDBI from the metrics
%% logger, which is the directly comparable number.
%% ===================================================================
\section{Experimental Results}
TODO

%% ===================================================================
%% V. CONCLUSION --- budget: ~250 w
%% ===================================================================
\section{Conclusion}
TODO

%% While drafting, include every entry in refs.bib so the reference list's
%% page cost stays visible. Once the body actually cites things, delete this
%% line -- refs.bib should end up being exactly the cited set.
\nocite{*}

\bibliographystyle{IEEEtran}
\bibliography{refs}

\end{document}
